\begin{statement}[C.6.T1; Proved, with exact local and descent computations]
At $t=32/5$ a global minimal model and the image of the section are
\[
 E:y^2=x^3-2x+1,\qquad P=(0,1),\qquad \Delta=80.
\]
The conductor is $40$, $E(\Q)=\Z/4\Z$ is generated by $P$,
and $\widehat h(P)=0$. The index of $\Z P$ in the quotient
modulo torsion is one. The complete bad-prime data are
\[
\begin{array}{c|c|c|c|c}
p&\operatorname{ord}_p\Delta&\text{type}&\text{component of }P&\lambda_p(P)\\\hline
2&4&\mathrm{III}&\text{unique nonzero class, order }2&\log2/12\\
5&1&\mathrm I_1&\text{identity}&\log5/12.
\end{array}
\]
All other finite local heights vanish. Consequently
\[
 \log N=\frac{\log10}{12},\qquad
 D_{\rm part}=\frac6\pi\log\frac{10^{1/12}}{2\sqrt3}.
\]
Numerically the two sides of the global height identity equal
\[
 0.191882091082837140334832621223697017300091790719064414669444\ldots.
\]
With the full real period $\Omega=5.9376498936895458116293755517\ldots$,
the central derivative is
\[
 \frac{L'(E,1)}\Omega=
 0.071331625531738774092617964261659879861112925567606085287841\ldots.
\]
The analytic rank is zero; a vanishing first derivative is therefore
not expected even though the point $P$ is torsion.
\end{statement}
\begin{proof}
The change of variables from the displayed family is
$x=(24/5)^2X$, $y=(24/5)^3Y$. Tate's algorithm gives the
listed minimal model, conductor and local data. Since $P$ has
order four, $3P=-P$. The local division identity consequently
gives
\[
 \lambda_p(P)=\left(\frac{v_p(\Delta)}{12}
                         -\frac{v_p(\psi_3(P))}{8}\right)\log p.
\]
Here $\psi_3(P)=-4$, proving both local coefficients exactly.
The exact descent output is $[0,0,0,[]]$, and the torsion
calculation gives $\Z/4\Z$. These computations, the direct
global equality and the analytic derivative are in
\path{code/Crank_torsion_local.gp} and
\path{code/Crank_torsion_final.gp}, with their outputs.
As a further check, the exact plus modular symbol from $0$ to
$\infty$ is $-1/4$ in the real-period lattice normalization;
the numerical ratio $L(E,1)/\Omega$ is $1/8$. The completed
functional equation then gives the checked expression
$(\gamma+\log(2\pi/\sqrt{40}))/8$ for the displayed derivative
ratio. The local-height proof does not use this $L$-value check.
\end{proof}
